% Record of expert interview / statement - Marta Parazzini (CNR-IEIIT), 22 June 2026.
% Co-written with Marta and lightly cleaned by Robin; preserves her quoted / faithfully paraphrased words.
% High-value market-discovery evidence for the IOF application (need, FTE waste, willingness to pay,
% adoption criteria, current licence spend). Signable.
% Compile: pdflatex parazzini_statement_22jun.tex
\documentclass[11pt]{article}
\usepackage[a4paper,margin=2.5cm]{geometry}
\usepackage{parskip}
\usepackage{graphicx}
\usepackage{enumitem}
\usepackage[T1]{fontenc}
\usepackage{lmodern}
\pagestyle{empty}

\begin{document}

% --- CNR-IEIIT letterhead placeholder ---------------------------------------
% \includegraphics[height=1.6cm]{cnr_ieiit_logo}\\[0.5em]
{\large\textbf{CNR-IEIIT}}\\
Istituto di Elettronica e di Ingegneria dell'Informazione e delle Telecomunicazioni\\
Consiglio Nazionale delle Ricerche, Milan, Italy\\[1.5em]
% ----------------------------------------------------------------------------

\begin{center}
\textbf{\large Statement on the need for fast numerical dosimetry}\\[0.3em]
\textit{Record of an interview in support of the IOF valorisation project (AEGIS)}
\end{center}

\bigskip
On 22 June 2026, I had a video conference call with Mr.\ Robin Wydaeghe. I was asked
questions in a non-leading, objective way. During this call, I am either quoted or faithfully
paraphrased in the following statements:

\medskip
\begin{itemize}[leftmargin=1.4em,itemsep=0.5em]

  \item There are 4 to 5 researchers working with FDTD in my research group. Two of them work
  in the high-frequency regime ($>6$\,GHz), three of them in the low-frequency regime for
  biomedical applications ($<6$\,GHz).

  \item With the advent of 6G comprising the new Frequency Band 3, I intend to hire more
  researchers working in the high-frequency regime in the coming years.

  \item Each of the employees in this category spends around 30--50\% of their time waiting
  for FDTD simulations to finish, which translates to 1.5 FTE in total (rising in the future).
  The rest of their time is spent on data analysis, paper writing, reading literature, and
  similar tasks.

  \item Configuring FDTD simulations takes around 10\% of a senior researcher's time and
  around 30\% of a junior researcher's time.

  \item The bottleneck in their research is always the mesh; it is our nightmare.

  \item Researchers perform high-frequency simulations in both the far-field (incident fields
  in an environment) and the near-field (incident from a device).

  \item Three things would help convince me to use a fast numerical dosimetry method, or to
  recommend it to others:
  \begin{enumerate}[leftmargin=1.6em,topsep=0.3em,itemsep=0.2em]
    \item Clear tutorials and an easy-to-use platform are extremely useful. This enables us to
    validate it ourselves while learning it. I rate the importance of this 5/5.
    \item A peer-reviewed paper, in the sense that it validates the method. I rate the
    importance of this 4/5.
    \item If I am performing standards-related work, it is useful if it is suggested as a
    viable method in the standards. I rate the importance of this 3/5.
  \end{enumerate}
  Point 1 is likely enough to convince me, with the others strengthening my decision.

  \item On the commercial value of the method: \textit{``If I really find a real improvement
  with respect to what I have now, I would be able to pay you.''}

  \item Currently, we own 4 licences necessary to run high-frequency exposure simulations,
  each at 3\,kEUR per year, and various other ones (such as those for phantoms) similarly
  priced.

  \item Zurich MedTech provides some licences for free as part of their University Research
  Agreement, but only when we use them for research purposes. We also pay for discounted CST
  licences.

\end{itemize}

\vspace{1.5em}
With best regards,

\vspace{2.2cm}

\rule{6cm}{0.4pt}\\
Dr.\ Marta Parazzini\\
Director of Research, CNR-IEIIT\\
Consiglio Nazionale delle Ricerche, Milan, Italy\\
marta.parazzini@cnr.it \hfill Date: [date]

\end{document}
